La diffusione dei modelli linguistici basati su Transformer rende centrale il costo dei loro pesi. Durante la generazione autoregressiva, i coefficienti delle proiezioni lineari devono essere disponibili a ogni nuovo token: il numero di parametri determina quindi una parte rilevante della memoria richiesta e del traffico fra memoria e unità di calcolo. Quando un modello supera la memoria disponibile, i suoi pesi devono essere trasferiti dal disco durante l'esecuzione, e il tempo di generazione diventa dominato dalla quantità di byte da leggere.

Questa tesi propone \textbf{STERA}, acronimo di \emph{Structured Training, Execution and Reconstruction Architecture}: una parametrizzazione strutturata dei pesi per le proiezioni lineari dei Transformer, usata direttamente durante il training, che riduce i byte necessari per memorizzare il modello e che si applica in inferenza senza ricostruire le matrici dense. Il termine \emph{Architecture} indica l'insieme delle procedure di training, esecuzione e ricostruzione associate alla parametrizzazione: STERA non introduce una nuova famiglia di architetture neurali, ma mantiene il Transformer autoregressivo e interviene solo sulla rappresentazione dei coefficienti delle sue proiezioni.

L'elemento centrale di STERA sono i blocchi \emph{RowMask23}. A differenza della compressione applicata a un modello già addestrato, il vincolo strutturale è imposto fin dall'inizio: l'optimizer aggiorna direttamente le coordinate dei blocchi e non esiste una matrice densa da approssimare. Il risparmio sulla parte strutturata delle matrici è quindi garantito per costruzione, mentre le componenti che completano i blocchi (scale, guadagni e termine low-rank) determinano il bilancio complessivo, che la tesi ricava in forma esplicita per qualunque dimensione del modello.

Il modello linguistico addestrato e valutato è \textbf{STERA 30M}: un Transformer a sei layer, con hidden size 768, 16 attention head e feed-forward di dimensione \num{3072}, le cui proiezioni sono parametrizzate con STERA. Il modello conta \num{30346182} parametri allenabili, che determinano \num{82721280} valori nella rappresentazione densa equivalente. È stato addestrato da zero sul corpus TinyStories per \num{122000} aggiornamenti, pari a \num{1529514422} token target.

\section*{Ipotesi}

La tesi verifica quattro ipotesi.
\begin{enumerate}
\item[\textbf{H1}] \textbf{Apprendimento.} Un Transformer le cui proiezioni sono parametrizzate con STERA apprende regolarità linguistiche che generalizzano a testi esclusi dal training.
\item[\textbf{H2}] \textbf{Rappresentazione.} La parametrizzazione riduce i byte necessari per rappresentare i pesi, rispetto alla propria espansione densa e a modelli densi del dominio; il risparmio è garantito dalla struttura dei blocchi e dipende in modo esplicito dalle altre componenti delle proiezioni.
\item[\textbf{H3}] \textbf{Inferenza.} Applicando le proiezioni direttamente alle coordinate, la generazione è più rapida che con le matrici dense equivalenti, a parità di pesi.
\item[\textbf{H4}] \textbf{Qualità narrativa.} Il modello genera racconti di qualità confrontabile con quella di modelli densi pubblici dello stesso dominio.
\end{enumerate}

\section*{Contributi}

\begin{itemize}
\item La definizione della parametrizzazione STERA e delle sue proprietà: la mappa dalle coordinate al blocco è lineare e iniettiva, il gradiente ha forma chiusa e il prodotto con un vettore si calcola con la trasformata di Walsh--Hadamard. Una formula esatta dà il risparmio di una proiezione e di un modello intero in funzione delle dimensioni e del rango.
\item Il training da zero di STERA 30M e il suo confronto con due checkpoint pubblici dello stesso dominio, TinyStories-8M e TinyStories-33M, sulla likelihood di racconti di prova e con un giudizio cieco delle continuazioni generate.
\item Due motori di esecuzione per gli stessi pesi, uno che applica le proiezioni direttamente alle coordinate e uno con le matrici dense equivalenti, con misure di memoria e tempi al variare della lunghezza dell'ingresso e un confronto con i modelli di riferimento sullo stesso motore.
\item Uno studio a posteriori del contributo delle componenti delle proiezioni (rango, scale e guadagni) e della quantizzazione a 8 bit del termine low-rank.
\end{itemize}

Gli esiti numerici sono riassunti nel Sommario e discussi nei Capitoli~7 e~8; il Capitolo~9 riporta l'esito di ciascuna ipotesi. Tutti i risultati si riferiscono al checkpoint finale di STERA 30M e sono riproducibili con il codice, le procedure e le impronte dei file disponibili nel repository del progetto.

\section*{Organizzazione del testo}

Il Capitolo~1 presenta il contesto e le motivazioni. Il Capitolo~2 richiama i fondamenti teorici: modelli autoregressivi, Transformer, training, generazione e trasferimento dei pesi fra memoria e disco. Il Capitolo~3 discute lo stato dell'arte sulla compressione dei modelli, sulle matrici strutturate e sull'esecuzione con memoria limitata. Il Capitolo~4 definisce la parametrizzazione STERA: i blocchi RowMask23 e le loro proprietà, le scale, i guadagni e il termine low-rank, e il bilancio complessivo dei parametri al crescere del modello. Il Capitolo~5 descrive STERA 30M e il percorso dei pesi dal training alla generazione. Il Capitolo~6 presenta i dati, il training e la metodologia di validazione, il Capitolo~7 i risultati su apprendimento e qualità dei racconti e il Capitolo~8 quelli su memoria e tempi di esecuzione. Il Capitolo~9 riassume l'esito delle ipotesi, i limiti del lavoro e gli sviluppi futuri.
